\documentclass[10pt]{article}

% Specify the margins and text size.
\setlength{\textwidth}{6.9in}
\setlength{\textheight}{9.5in}
\setlength{\oddsidemargin}{0pt}
\setlength{\evensidemargin}{0pt}
\setlength{\topmargin}{0pt}
\setlength{\hoffset}{-.15in}
\setlength{\voffset}{-1in}

\setlength{\parskip}{5pt}
\setlength{\parindent}{0pt}

% Load some fonts and symbol packages
\usepackage{latexsym}
\usepackage{pifont}       % contains 'star' symbol for counterinsurgency handbook title
\usepackage{yfonts} 
\usepackage{amsmath}
\usepackage{amsfonts}

\usepackage{graphicx}     % actually, this is loaded by pstricks
\usepackage[T1]{fontenc}
\usepackage{ifthen}
\usepackage{pstricks,pst-grad,pst-text,pst-node,multido,pst-plot,pst-func,calc,pst-3dplot}
%\usepackage[all]{xy}
%\usepackage{animate}

% The hyperref package inserts links into the pdf.
\definecolor{MyLinkColor}{rgb}{.1,.2,1}
\definecolor{MyCiteColor}{rgb}{.1,1,.2}
\definecolor{MyURLColor}{rgb}{.4,.4,.4}
\usepackage[backref=true,pagebackref=false,hyperindex,colorlinks=true,
  linkcolor=MyLinkColor,urlcolor=MyURLColor]{hyperref}


% The tweaklist package is something I found on the web.  It provides a simple interface
% for making changes to spacing used in the itemize and enumerate environments.  Comment
% this out if you don't care to use tweaklist.
\usepackage{tweaklist}
\renewcommand{\itemhook}{\setlength{\parskip}{2pt}\setlength{\parsep}%
{1pt}\setlength{\topsep}{0pt}\setlength{\itemsep}{0pt}}

\newcommand{\U}{\underline{\hspace{5pt}}}

\usepackage{listings}
\newcommand{\Z}{\hphantom{0}}

\begin{document}
\pagestyle{empty}
\lstset{language=R, showspaces=false, showstringspaces=false}
\href{http://www.su.edu}{\includegraphics[height=1.75cm]{sulogo.eps}}
\vspace{-1.69cm}

{{\ }\hfill\small
\begin{tabular}{cl}
& Math 207\\
& Introduction to Statistics\\
\end{tabular}
}
\setlength{\baselineskip}{1.05\baselineskip}

\begin{center}
\textbf{\large  Review Exercises}
\end{center}


1. Four hundred draws are made at random with replacement from a box
of numbered tickets;  150 are positive.  Someone tells you that
50\% of the tickets in the box show positive numbers.  Do you believe it?  
Answer yes or no, and explain using a $P$-value.  Then find a 95\% confidence interval
for the percent of positive tickets in the box.
\vspace{1.8in}


2. With a perfectly balanced roulette wheel, in the long run, 
red numbers should turn up 18 times in 38.
To test its wheel, one casino records the results of 400 plays, finding 200 red
numbers.
%To test its wheel, one casino records the results of 3,800 plays, finding 1,890 red
%numbers.  
Is that too many reds?  Or chance variation?
Formulate the null and alternative hypotheses as 
statements about a box model.
Compute $z$ and $P$.
\vspace{1.5in}

3. Tuddenham and Snyder obtained the following results for 66 California children
at ages 6 and 18.  \vspace{-6pt}
\begin{center}
\begin{tabular}{rclrcl}
average height at age 6 & $\approx$ & 3 feet 10 inches, & SD &$\approx$ & 1.7 inches\\
average height at age 18 & $\approx$ & 5 feet 10 inches, & SD &$\approx$ & 2.5 inches
  \hspace{10pt} $r\approx 0.80$\\
\end{tabular}\vspace{-6pt}
\end{center}
The scatter diagram is football-shaped.  

(a) Find the regression line and RMS error for predicting height at 18 from height at 6.
\vspace{1in}

(b) Find the regression line and RMS error for predicting height at 6 from height at 18.
\vspace{1in}

(c) Of children who were 4 feet at age 6, 95\% of those had heights
at age 18 in what range?

\vfill
\eject
\newcommand{\BL}{\underbar{\hspace{.5in}}}



%\#1 from page 455.
4.
Laser altimeters can measure elevation to within a few inches, without bias,
and with no trend or pattern to the measurements.
As part of an experiment, 25 readings were made on the elevation of a mountain peak.
These averaged out to 81,411 inches, and their SD was 30 inches. 

\hspace{20pt} a) The elevation of the mountain peak is estimated 
as \BL; this estimate is likely to be off by\vspace{-4pt}

\hspace{20pt}\hphantom{f) } \BL\ or so.
\bigskip
\bigskip

\hspace{20pt} b) (T/F) 81,411$\pm$ 12 inches is a 95\%-confidence interval for the elevation
of the mountain peak. 
\bigskip
\bigskip

\hspace{20pt} c) (T/F) 81,411$\pm$ 12 inches is a 95\%-confidence interval for 
average of the 25 readings.

\hspace{20pt} d) (T/F) There is about a 95\% chance that the next reading will be in 
the range 81,411 $\pm$ 12 inches.

\hspace{20pt} e) (T/F) About 95\% of the readings were in the range 81,411 $\pm$ 12 inches.

\hspace{20pt} f) If another 25 readings are made, there is about a 95\% chance that their
average will be in the range\vspace{-4pt}

\hspace{20pt}\hphantom{f) } 81,411 $\pm$ 12 inches.
\bigskip\bigskip

5. A coin is to be flipped 100 times. What is the chance of 
getting exactly 97 heads?  Simplify your answer as much as possible.
What is the chance of getting at least 97 heads?
\vspace{1.5in}


6.  A spectrophotometer is correctly calibrated if its average measurement is~70~ppm.
Three test measurements are taken.  The values are 71, 83, and 74.  Use  a
$t$-test to see whether the instrument is correctly calibrated.  
State the null and alternative hypothesis. Calculate the $P$ value.
(Hint:  $\sqrt{39}\approx 6$.)  Use the table and plot below to calculate
areas under  $t$-distributions.
\smallskip

\begin{tabular}{ccc}
\begin{tabular}{|c|cccc|}\hline
Degrees\vphantom{\LARGE Y} & &&&\\
of      & &\rput(20pt, 10pt){$t$}&&\\
Freedom & $-2.0$ & $-1.5$ & $-1.0$ & $-0.5$\\[3pt]\hline
1       & $0.148$ & $0.187$ & $0.250$ & $0.352$\vphantom{\LARGE Y}\\
2       & $0.092$ & $0.136$ & $0.211$ & $0.333$\\
3       & $0.070$ & $0.115$ & $0.196$ & $0.326$\\\hline
\end{tabular}
&\hspace{1in} &
\definecolor{lg}{rgb}{.8,.8,.8}
\begin{pspicture}(-2.5,0.5)(2.5,1.5)
\psframe[linewidth=0.02](-3.2,-0.88)(3.2,2.05)
\rput(0,1.7){$t$ distribution with $2$ degrees of freedom}
\psset{yunit=4}
\pscustom[linestyle=none,fillstyle=solid,fillcolor=lg]{%
\psTDist*[linewidth=0,linecolor=blue,nue=2]{-2.25}{-0.5}
\psline(-0.5,0.2963)(-0.5,0)(-2.25,0)(-2.25,0.0533)
}
\psTDist[linecolor=blue,nue=2]{-2.25}{2.25}
\psline{<->}(-2.5,0)(2.5,0)
\psline(-0.5,0)(-0.5,-0.05)\rput[t](-0.5,-0.08){\footnotesize$-0.5\hphantom{-}$}
\psline(-2,0)(-2,-0.05)    \rput[t](-2,-0.08){\footnotesize$-2\hphantom{-}$}
\psline(-1,0)(-1,-0.05)   % \rput[tr](-1,-0.08){\footnotesize$-1$}
\psline(0,0)(0,-0.05)      \rput[t]( 0,-0.08){\footnotesize$0$}
\psline(1,0)(1,-0.05)      \rput[t]( 1,-0.08){\footnotesize$1$}
\psline(2,0)(2,-0.05)      \rput[t]( 2,-0.08){\footnotesize$2$}
\psbezier{->}(-6.13,-0.05)(-4,0.3)(-2,0.5)(-1,0.1)
\end{pspicture}
\end{tabular}
\vspace{2in}
\vfill
\eject

%\#4 from page 506.
7. A study reports that freshmen at public universities work 10.2 hours a week for pay,
on average, and the SD is 8.5 hours;
at private universities, the average is 8.1 hours and the SD is 6.9 hours.
Assume these data are based on two independent simple random samples, each of
size 1,000.  Is the difference between the
averages due to chance?  If not, what else might explain it?
\vspace{2in}

8.
Suppose that you take a large number of independent random samples from a box of numbered
tickets.
Fill in the blanks.

\begin{itemize}
\item[a)] The histogram of the averages of the samples will
  will look approximately like a \underbar{\hspace{1.2in}}.\vspace{.2in}
\item[b)] The histogram of the sums of the  samples will
  will look approximately like a \underbar{\hspace{1.2in}}.\vspace{.2in}
  
\item[c)] Suppose that the tickets are all 0s and 1s.  The histogram of the
  percent of 1s in the  samples will
  will look approximately like a \underbar{\hspace{1.2in}}.
\end{itemize}
\medskip

9.  One hundred draws will be made at random with replacement from
the box
\begin{center}
\begin{tabular}{|cccccc|}
0 &0 & 2 & 2 & 2 & 6\\\hline
\end{tabular}
\end{center}

\hspace{10pt} a) What is the smallest that the sum can be?
\vspace{.5in}

\hspace{10pt} b) What is the largest that the sum can be?
\vspace{.5in}

\hspace{10pt} c) The sum of the draws will be around \underline{\hspace{1in}},
give or take \underline{\hspace{1in}} or so.  %Justify your answers.
\vspace{1in}

\hspace{10pt} d) The chance that the sum will be less than 175 is around
\underline{\hspace{1in}} \%.  Justify your answer.


\vfill
\eject

10.  An experiment consists of rolling a die 900 times and recording how often
the die shows 1 or 6.

\hspace{10pt} a) Write down a box model to represent the experiment.
\vspace{.5in}

\hspace{10pt} b) Compute the mean and SD of the box.
\vspace{1in}

\hspace{10pt} c) What is the expected number of times the die will show a 1 or 6?
\vspace{1in}

\hspace{10pt} d) What expected percentage of times the die will show a 1 or 6?
\vspace{1in}

\hspace{10pt} e) The number of times that the die will show a 1 or a six 
will be around \underline{\hspace{1in}},\vspace{10pt}

\hspace{10pt}\hphantom{e) }give or take \underline{\hspace{1in}} or so.  %Justify your answers.
\vspace{1.5in}

\hspace{10pt} f) The chance that the percent of times will be between 31\% and 35\% is about
\underline{\hspace{1in}} \%.  
\vfill
\eject

\begin{center}
\textbf{\large Solutions}
\end{center}

1. The observed fraction is $150/400=37.5\%$.  
The null hypothesis is that 50\% of the tickets are positive and that the difference between
the expected and observed values is due to chance.  If the null hypothesis is true, then
the SD of the box is $\sqrt{0.5\times 0.5}=0.5$.  So, the expected value of the percent of
positives is 50\% and the SE is $0.5/\sqrt{400}=2.5\%$.  That gives a $z$-score of 
$(37.5-50)/2.5=-5$.  The $p$-value is $\mbox{pnorm}(-5)\approx 0.000003\%$.  That supports the 
alternative hypothesis.  So, I don't believe the 50\% claim. 

To get a 95\% confidence interval on the true percent of positives, use $37.5\%$ to compute 
$\mbox{SD}=\sqrt{0.375\times (1-0.375)}=0.47$ and $\mbox{SE}=0.47/\sqrt{400}=2.35\%$.  So,
the confidence interval is $37.5\% \pm 2\times 2.35\% = 37.5\% \pm 4.7\%$.
\medskip

2. The observed percent of reds is $200/400=50\%$.  The box model has 18 ones (representing the reds)
and 20 zeros (representing the other possibilities).  The average of the box is $18/38 = 47.37\%$
and its SD is $\sqrt{.4737\times (1-.4737)}\approx 0.5$.  The expected value for the percent of
reds is 47.37\% and the SE for 400 spins is $0.5/\sqrt{400}=2.5\%$.  The null hypothesis is that
the observed difference (50\% vs 47.37\%) is due to chance.  The alternative is that the 
difference is not due to chance.  The $z$-score is $(50-47.37)/2.5 = 1.052$ and the
$p$-value is $1 - \mbox{pnorm}(1.052)=14.6\%$.  This supports the null hypothesis.  The roulette
wheel seems to be working correctly.
\medskip

3. a) Use the age 18 data for $y$ and the age 6 data for $x$. 
The regression line is $y-70 = 0.8\times (2.5/1.7)\times (x - 46)$
That is $ y - 70 = 1.18\,(x - 46)$.  The RMS error is $\mbox{SD}_y\sqrt{1-r^2}=
2.5\,\sqrt{1- 0.8^2} = 1.5$ inches.

b) Use the age 6 data for $y$ and the age 18 data for $x$. 
The regression line is $y-46 = 0.8\times (1.7/2.5)\times (x - 70)$
That is $ y - 46 = 0.544\,(x - 70)$.  The RMS error is $\mbox{SD}_y\sqrt{1-r^2}=
1.7\,\sqrt{1- 0.8^2} = 1.02$ inches.

c) Use the formulas from a).  If $x=48$ inches, then the predicted $y$ is
72.36 inches.  95\% of the data should fall in the range $72.36\pm 2\times 1.5$ inches.
\medskip

4. a) First blank is the observed value:  81,411 inches.  The second blank is the 
standard error:  $30/\sqrt{25}=6$ inches.

b) True.  observed value $\pm$ 2 standard errors.

c) False.  If we have 25 measurements, we can take their average.  There is no $\pm$.

d)  False.  The SD for individual readings is 30 inches. So, this answer should be $\pm 60$ inches.

e) False.  Same reason as d).
 
f)  False.  We are 95\% sure that the true value will be in a confidence interval.  We don't
have that confidence that the next average reading will be in our previous confidence interval.
\medskip

5. For exactly 97 heads, 
$\displaystyle \frac{100!}{97!\, 3!}\,\left(\frac{1}{2}\right)^{97}\,\left(\frac{1}{2}\right)^{3}
=\frac{100\cdot 99\cdot 98}{3\cdot 2}\left(\frac{1}{2}\right)^{100}
=161,\!700\left(\frac{1}{2}\right)^{100}$  It's a really small number.  In R you can compute it
with $\mbox{dbinom}(97, 100, 0.5)$.

For 97 or more heads, 
$\displaystyle \frac{100!}{97!\, 3!}\,\left(\frac{1}{2}\right)^{97}\left(\frac{1}{2}\right)^{3}
\; +\;\frac{100!}{98!\, 2!}\,\left(\frac{1}{2}\right)^{98}\,\left(\frac{1}{2}\right)^{2} \;+\;
\frac{100!}{99!\, 1!}\,\left(\frac{1}{2}\right)^{99}\,\left(\frac{1}{2}\right)^{1}\; +\;
\frac{100!}{100!\, 0!}\,\left(\frac{1}{2}\right)^{100}\,\left(\frac{1}{2}\right)^{0}$
which is $(161,\!700 + 4950 + 100 + 1)\,(1/2)^{100}$.
\medskip

6.  The observed average is $(71 + 83 + 74)/3 = 76$.  The observed sd is
$\sqrt{\frac{(71-76)^2 + (83-76)^2 + (74-76)^2}{2}}=6.2$.  The standard error (SE)
is $6.2/\sqrt{3}=3.6$.  
The null hypothesis is that
the difference between the observed (76) and expected (70) averages is due to chance.
The alternative is that the difference is not due to chance.  The test statistic
is $(76-70)/3.6 = 1.7$.  The $p$-value is $1 - \mbox{pt}(1.7, \, \mbox{df}=2)=11.6\%$.
This supports the null hypothesis.
\medskip

7. This is a two-sample $z$-test.  Box $A$ contains work hours for freshmen at public universities.
Box $B$ contains work hours for freshmen at private universities.  The sample from box $A$ had
average 10.2 hours and SD of 8.5 hours. The sample from box $B$ had average 8.1 hours and SD of 6.9
hours.  The SEs for the two boxes are estimated to be $8.5/\sqrt{1000}=0.27$ hours and
$6.9/\sqrt{1000}=0.22$ hours.  The SE for the difference is $\sqrt{0.27^2 + 0.22^2} = 0.35$ hours.
This gives a $z$-score of $(10.2 - 8.1)/0.35 = 6$.  The $p$-value is $1-\mbox{pnorm}(6)\approx 0$.
So, the difference is probably not due to chance.
\vfill
\eject

8. By the Central Limit Theorem, the answer for each blank is \textit{normal curve}.
\medskip

9. a) $100\times 0 = 0$.

b) $100\times 6 = 600$.

c) The mean and SD of the box are both $2$.  The expected value of the sum is 
$100\times 2 = 200$.  The standard error for the sum is $\sqrt{100}\times 2 = 20$.

d) Convert 175 to a  $z$-score:  $(175-200)/20 = -1.25$.  So the answer is 
$\mbox{pnorm}(-1.25)=10.6\%$.
\medskip

10. a) The box has 4 zeros (representing 2, 3, 4 and 5) and 2 ones (representing 1 and 6).

b) The mean is $1/3$ and the SD is $\sqrt{2}/3 = 0.47$.  

c) The expected value of the sum is $900\times (1/3) = 300$.  

d) The expected value of the percent is 33.3\%.

e)  The SE for the sum is $\sqrt{900}\, \sqrt{2}/3 = 30\,\sqrt{2}/3 = 10\sqrt{2}=14.1$.
So, the answer is $300\pm 14.1$.

f) The SE for the percent is $0.47/\sqrt{900} = 1.6\%$.
Convert 31\% and 35\% to $z$-scores:  
$(31-33.3)/1.6=-1.4$ and $(35-33.3)/1.6=1.1$.  So, the answer is
$\mbox{pnorm}(1.1) - \mbox{pnorm}(-1.4)=78.4\%$.
\vfill
\eject
\end{document}

2.  The spinner for a game-of-chance is tested to see if it is balanced
by spinning it $50$ times. 
Use the data in the table below to test the null hypothesis that that spinner is
fair (i.e.,~that the differences between the observed and expected frequencies are
due to chance). Compute the $P$-value and discuss its significance.
\begin{center}
\begin{tabular}{ccc}
\begin{pspicture}(-2,-2)(2,2)
\pswedge*[linecolor=lightgray](0,0){2}{0}{72}
\pswedge[hatchcolor=lightgray,fillstyle=crosshatch](0,0){2}{72}{144}
\pswedge*[linecolor=white](0,0){2}{144}{216}
\pswedge[hatchcolor=white,fillstyle=crosshatch*,fillcolor=lightgray](0,0){2}{216}{288}
\pswedge[hatchcolor=lightgray,fillstyle=crosshatch](0,0){2}{288}{360}
\pscircle(0,0){2}
\SpecialCoor
\rput{126}(1.6;36){\textbf{A}}
\rput{198}(1.6;108){\textbf{B}}
\rput{270}(1.6;180){\textbf{C}}
\rput{-18}(1.6;252){\textbf{D}}
\rput{54}(1.6;-36){\textbf{E}}
\psset{linewidth=3pt}
\psline{*->}(1.0;-20)(1.4;160)
\pscircle[fillstyle=solid,linewidth=0.04,fillcolor=white](0,0){0.15}
\end{pspicture}
&
\hspace{1in}
&
\begin{pspicture}(-2,-2)(2,2)
\rput(0,0){\begin{tabular}{|c|c|}\hline
        & observed\vphantom{\Large Y}\\
  value & frequency\\[2pt]\hline
    A   &  \Z8\vphantom{\Large Y}\\
    B   &  12\\
    C   & \Z7\\
    D   & \Z9\\
    E   & 14\\\hline
\end{tabular}}
\end{pspicture}
\end{tabular}
\end{center}
\vfill
{\small
\begin{tabular}{|c|rrrrrrr|}\hline
df & 0.5 & 1.0 & 1.5 & 2.0\rput(-19pt,20pt){Table entries are area to the right of the given $\chi^2$ value} & 2.5 & 3.0 & 3.5\vphantom{\Large Y}\\\hline
1 &  0.48 & 0.32 & 0.22 & 0.16 & 0.11 & 0.08 & 0.06\vphantom{\Large Y}\\
2 &  0.78 & 0.61 & 0.47 & 0.37 & 0.29 & 0.22 & 0.17\\
3 &  0.92 & 0.80 & 0.68 & 0.57 & 0.48 & 0.39 & 0.32\\
4  & 0.97 & 0.91 & 0.83 & 0.74 & 0.64 & 0.56 & 0.48\\
5  & 0.99 & 0.96 & 0.91 & 0.85 & 0.78 & 0.70 & 0.62\\
6  & 1.00 & 0.99 & 0.96 & 0.92 & 0.87 & 0.81 & 0.74\\\hline
\end{tabular}}
\eject

\end{document}


\#3 from page 444.  (Hypothetical)  The British Imperial Yard is sent to Paris for calibration
against The Meter.  The length is determined 100 times.  This sequence of
measurements averages out to 91.4402 cm, and the SD is 800 microns (a \textit{micron}
is a millionth of a meter).

\hspace{20pt} a) Is a single reading off by around 80 microns, or 800 microns?
\bigskip

\hspace{20pt} b) Is the average of all 100 readings off by around 80 microns, or 800 microns?
\bigskip

\hspace{20pt} c) Find a 95\%-confidence interval for the exact length of the
Imperial Yard.
\vspace{.5in}

\#5 from page 253.
The Bureau is about to weigh a one-kilogram checkweight 100 times, and take
the average of the measurements.  They are willing to assume the
Gauss model, with no bias, and on the basis of past experience they 
estimate the SD of the error box to be 50 micrograms.

\hspace{20pt} a) The average of all 100 measurements is likely to be off
the exact weight by \BL\ or so.
\bigskip

\hspace{20pt} b) The SD of all 100 measurements is likely to be around \BL.
\bigskip
\bigskip

\hspace{20pt} c) Estimate the probability that the average of all 100 measurements
will be within 10 micrograms of the exact weight.
\vspace{.7in}
